\section*{The Propers: Detailed Commentary}

\subsection*{Urgency and unpriced covenant\enspace\elementref{Entrance, First Reading}}
Psalm 69 gives the entrance no ornamental calm: the worshipper asks God to hasten. Isaiah 55 answers need in a different register, inviting the thirsty and penniless to water, wine, milk, rich food, attentive hearing, and life. The invitation concludes Isaiah 40--55 and renews the “sure mercies of David”; it is covenant speech to displaced people, not a timeless advertisement for comfort.

Augustine repeatedly receives the prophetic “without money” as a grammar of grace: God's gift cannot be bought by prior merit. That reading does not erase the command to listen or the social obligations surrounding food. It identifies the giver's initiative. The Entrance and Isaiah can therefore be heard together as petition and answer without claiming that the Missal antiphon was historically selected for this Gospel cycle.

\subsection*{The open hand and the whole psalm\enspace\elementref{Responsorial Psalm}}
Psalm 145's appointed verses praise compassion, fidelity, justice, food in season, an open hand, and nearness to truthful prayer. The full psalm also names prisoners, the bowed down, strangers, orphans, and widows. Satisfaction is not abstract plenitude but the kingship of God turned toward dependent creatures and vulnerable persons.

Augustine's exposition treats divine praise as the creature's fitting answer to a King whose reign does not pass away. The liturgical response focuses the assembly on God's opened hand; the surrounding psalm prevents that phrase from shrinking into private consumption. Praise trains attention to those whom divine kingship protects.

\subsection*{Compassion, inadequacy, and mediated abundance\enspace\elementref{Gospel}}
Matthew places the feeding immediately after Herod's banquet and John the Baptist's death. Jesus withdraws, but the crowd follows; his first response is compassion and healing. When evening comes, the disciples propose dismissal. Jesus' answer does not deny their inventory of five loaves and two fish. It changes the question from possession to availability: what they have is brought to him.

Chrysostom's homily on the feeding stresses that Christ could have created food directly but leads the disciples through acknowledgment, obedience, and distribution. The blessing identifies the Father as giver; the fragments show abundance without turning waste into spectacle. Twelve baskets mark a surplus the disciples themselves must gather.

The narrative carries Eucharistic resonance through taking, blessing, breaking, and giving, and the Church has long received it typologically. Yet the historical meal remains bodily mercy, and Matthew's wording does not license a simplistic identity between miracle and sacrament. The crowd's hunger must not be spiritualized away; Eucharistic reading should deepen, not replace, concern for food insecurity.

\subsection*{A word that feeds\enspace\elementref{Acclamation}}
Matthew 4:4 cites Deuteronomy after Jesus' fast. The saying rejects the temptation to secure identity through autonomous display, not bread's goodness. Positioned before a feeding Gospel, it prevents two opposite reductions: people cannot live by calories alone, and fidelity to God's word cannot become an excuse to dismiss bodily need.

\subsection*{Tribulation cannot become separation\enspace\elementref{Second Reading}}
Romans 8 concludes with a sequence of real threats: anguish, persecution, famine, nakedness, peril, sword, death, life, powers, height, depth. Paul does not call them illusions. He denies their capacity to separate believers from God's love in Christ.

Chrysostom reads the catalogue as confidence born from the love already shown in Christ rather than from an easy life. Aquinas distinguishes the evils suffered from the divine good that sustains the person through them. This is hope without moral laundering: suffering may become the field of fidelity, but evil is not made good and perpetrators are not excused. The passage cannot demand passive exposure to abuse or displacement of practical protection.

Romans is semi-continuous in Ordinary Time. Its relation to the feeding narrative is canonical and pastoral rather than an officially correlated selection. It can strengthen trust when resources fail, but it should not be made into proof that every shortage will be miraculously reversed.

\subsection*{Restored creatures become an offering\enspace\elementref{Collect, Prayer over Offerings}}
The Collect addresses God as creator and governor, asking both restoration and preservation. The offering prayer then asks that gifts be sanctified and accepted while the worshippers themselves become an enduring gift. The movement is not from worthless creation to escape, but from damaged creatures through divine restoration toward stable self-offering.

This guards application of the feeding: resources are not despised because small, and people are not reduced to instruments of distribution. God receives gifts and persons, sanctifies rather than annihilates them, and orders service toward communion.

\subsection*{Bread, refreshment, and redemption\enspace\elementref{Communion, Prayer after Communion}}
Wisdom 16 rereads manna as heavenly nourishment suited to desire; John 6 identifies Christ himself as the bread of life. These are genuine alternatives, not two halves of one enacted antiphon. Wisdom emphasizes the giver's pedagogy in Israel's history. John emphasizes coming and believing in the Son.

The final prayer speaks after reception: heavenly refreshment is already gift, while continuing protection and worthiness for eternal redemption remain petitions. Sacramental reception is therefore neither bare remembrance nor automatic completion. It begins and sustains a life dependent upon grace.
